\chapter{АШИГЛАХ ТЕХНОЛОГИ}

Энэхүү бүлэгт төслийн техникийн шийдлүүдийг сонгосон шалтгааныг тайлбарлана. Технологийн стек сонгохдоо дараах шалгуурыг баримталсан: нэг хэл дээр бүх давхаргыг хэрэгжүүлэх боломжтой байх, орчин үеийн AI tool-calling стандарттай нийцэх, мөн семестрийн хугацаанд практик хөгжүүлэлт хийх боломжтой байх.

\section{Хөгжүүлэлтийн хэл}

Төслийн бүх давхаргад TypeScript хэлийг ашиглах шийдлийг гаргасан. TypeScript нь JavaScript-ийн дээд багц бөгөөд статик төрлийн системийг (static type system) нэмж өгдөг. Төслийн хувьд TypeScript-ийн давуу тал нь Chrome extension болон MCP интеграцын хооронд ижил төрлийн тодорхойлолт (type definition) ашиглах боломжтой байдагт оршино. Жишээлбэл, MCP хэрэгслийн оролт гаралтын бүтцийг нэг удаа тодорхойлоод, AI сервер болон Chrome extension хоёуланд нь ашиглаж болно.

TypeScript нь хөгжүүлэлтийн явцад алдааг эрт илрүүлж, IDE-ийн автомат гүйцээлт (autocomplete) болон рефакторинг хийхэд тустай. Энэ нь ялангуяа олон модулиас бүрдсэн monorepo бүтэцтэй төсөлд чухал ач холбогдолтой юм.

\section{Model Context Protocol (MCP)}

Model Context Protocol SDK нь MCP стандартын албан ёсны TypeScript хэрэгжүүлэлт юм \cite{mcp-spec}. Энэхүү SDK нь MCP серверийг үүсгэж, tool, resource, prompt-уудыг бүртгэх боломж олгодог. Хэрэгслийн оролт гаралтын бүтцийг Zod schema ашиглан тодорхойлдог бөгөөд энэ нь runtime validation болон TypeScript type inference хоёуланг хангадаг. Тээврийн давхаргын хувьд stdio, HTTP, WebSocket зэрэг сонголтуудыг дэмждэг.

Одоогийн хэрэгжилтэд stdio transport ашиглан локал MCP сервер ажиллуулж байна. MCP серверийн хэрэгжилт нь төслийн гол техникийн бүрэлдэхүүн хэсэг бөгөөд SISI порталын API-тай холбогдох найман read-only хэрэгсэл тодорхойлогдсон.

\section{Chrome Extension}

Chrome extension хөгжүүлэхэд WXT framework ашиглаж байна \cite{wxt}. WXT нь орчин үеийн extension хөгжүүлэлтийн фреймворк бөгөөд уламжлалт extension хөгжүүлэлттэй харьцуулахад хэд хэдэн давуу талтай. Хөгжүүлэлтийн явцад Hot Module Replacement (HMR) ашиглан extension-ийг дахин ачаалахгүйгээр өөрчлөлтийг харах боломжтой. Extension-ийн бүх хэсэгт TypeScript ашиглах боломжтой бөгөөд шинэ Manifest V3 стандартыг бүрэн дэмждэг.

Extension-ийн хэрэглэгчийн интерфейсийг React ашиглан бүтээж байна. React нь компонент дээр суурилсан архитектур, виртуал DOM, өргөн хүрээний экосистем зэргээрээ extension UI хөгжүүлэхэд тохиромжтой. Popup доторх чат интерфейс нь хэрэглэгчийн үндсэн харилцааны цонх болж ажилладаг.

\section{AI интеграц}

AI загвартай ажиллахад Vercel AI SDK ашиглахаар төлөвлөж байна \cite{vercel-ai-sdk}. Энэ SDK нь хиймэл оюуны загваруудтай ажиллах TypeScript сан бөгөөд хэд хэдэн чухал боломжуудыг олгодог. Server-Sent Events ашиглан хэрэглэгчдэд бодит цагийн streaming хариу үзүүлэх боломжтой бөгөөд энэ нь урт хариулт үүсгэх үед хэрэглэгчийн туршлагыг сайжруулдаг. AI загварт хэрэгсэл тодорхойлж, загварын дуудлагыг автоматаар боловсруулах tool calling дэмжлэгтэй бөгөөд энэ нь MCP хэрэгслүүдтэй холбогдоход чухал. OpenAI, Anthropic, Google Gemini зэрэг үйлчилгээ үзүүлэгчдийг ижил интерфейсээр ашиглах боломжтой учраас ирээдүйд загвар солих уян хатан байдлыг хангадаг.

Прототип хөгжүүлэлтийн шатанд Google-ийн Gemini загварыг ашиглахаар төлөвлөж байна. Flash загвар нь хурдан хариу өгдөг бөгөөд прототип турших, хөгжүүлэлт хийхэд тохиромжтой. Gemini нь function calling-ийг сайн дэмждэг бөгөөд MCP хэрэгслүүдийг дуудахад ашиглах боломжтой. API дуудлагын өртөг харьцангуй бага учраас хөгжүүлэлтийн шатанд ашиглахад тохиромжтой. Эцсийн бүтээгдэхүүнд загварыг солих боломжтой бөгөөд Vercel AI SDK нь үүнийг хялбар болгодог.

\section{Monorepo бүтэц}

Төслийн кодын санг monorepo хэлбэрээр зохион байгуулсан. Monorepo нь олон апп, package-ийг нэг Git repository-д хадгалах арга бөгөөд код хуваалцах, хамаарлыг удирдахад тохиромжтой. Turborepo ашиглан monorepo-г удирдаж байгаа бөгөөд энэ нь build, test зэрэг даалгавруудын үр дүнг кэшлэж, хамааралгүй даалгавруудыг зэрэгцүүлэн ажиллуулдаг.

Одоогийн monorepo бүтэц нь apps болон packages гэсэн хоёр үндсэн хавтастай. Apps хавтаст extension (Chrome extension), server (AI сервер), web (туршилтын вэб интерфейс) багтдаг. Packages хавтаст mcp (MCP сервер), ui (дундын UI компонентууд), env (орчны тохиргоо) зэрэг дундын модулиуд багтдаг.

\section{Технологийн стекийн нэгдсэн хүснэгт}

Хүснэгт~\ref{tab:stack}--д төслийн технологийн стекийг нэгтгэн үзүүлэв.

\begin{table}[H]
  \centering
  \caption{Технологийн стек ба сонгосон шалтгаан}
  \label{tab:stack}
  \begin{tabularx}{\textwidth}{>{\raggedright\arraybackslash}p{3.3cm} >{\raggedright\arraybackslash}p{3.2cm} X}
    \toprule
    Ангилал & Технологи & Сонгосон шалтгаан \\
    \midrule
    Хэл & TypeScript & Extension, MCP интеграцын хооронд ижил type system ашиглах боломжтой \\
    MCP & Model Context Protocol SDK & Стандартчилсан tool, resource тодорхойлолт, Zod schema validation \\
    Extension & WXT + React & HMR, TypeScript дэмжлэг, Manifest V3, компонент архитектур \\
    AI давхарга & Vercel AI SDK & Streaming, tool calling, олон загварын дэмжлэг \\
    Загвар & Gemini & Хурдан хариу, tool calling дэмжлэг, бага өртөг \\
    Monorepo & Turborepo & Task caching, parallel execution, dependency graph \\
    \bottomrule
  \end{tabularx}
\end{table}

\section{Цаашид нэмэгдэх технологи}

Дараагийн шатанд хэд хэдэн технологийн хэсэг нэмэгдэх төлөвтэй байна. Гадаад MCP үйлчилгээнүүдтэй (Google Calendar, Notion гэх мэт) холбогдоход OAuth 2.0 зөвшөөрлийн урсгал шаардлагатай болно. Багшийн порталын Excel файлаас дүн оруулах workflow-д xlsx сан ашиглагдана. Mutation үйлдлүүдийг баталгаажуулах, хэрэглэгчээс зөвшөөрөл авах approval layer хэрэгжүүлэх шаардлагатай. Эдгээр өргөтгөл нь одоогийн суурь архитектуртай зөрчилдөхгүй, харин модульчлагдсан хэлбэрээр дээр нь нэмэгдэхээр төлөвлөсөн.
